%% ====================================================================
%%  How to Use This Guide  (unnumbered preface)
%% ====================================================================
\chapter*{How to Use This Guide}
\addcontentsline{toc}{chapter}{How to Use This Guide}
\markboth{How to Use This Guide}{}

This document is a \emph{concept-first companion} to the thesis \emph{``Lightweight
Knowledge-Distilled Deep Learning for WiFi Indoor Localization.''} The thesis reports \emph{what} we
did and \emph{what} we found; this guide explains \emph{why} every concept works and \emph{why} every
model behaved the way it did, starting from the foundations. It is written so that each of the three
authors can read any part independently and walk into the defense able to explain --- not just recite
--- the work. It complements the four system-reference documents in \code{docs/}
(\code{01\_dloc\_architecture}, \code{02\_dataset}, \code{03\_model\_architectures},
\code{04\_experiments}); where those say ``here is exactly what our system is,'' this guide teaches
the underlying ideas.

\section*{The thesis in one page}

\textbf{The problem.} GPS does not work indoors, but WiFi access points (APs) are everywhere and
already measure \emph{Channel State Information} (CSI). DLoc~--- a 2020 system from UC San Diego ---
turns CSI into an image and learns to ``translate'' it into a heatmap whose bright spot is the
person's location, reaching sub-meter accuracy. The catch: DLoc has \textbf{16.5 million parameters},
far too heavy for the edge devices (gateways, phones, the APs themselves) where indoor localization
is actually useful.

\textbf{The question.} Is that 16.5\,M cost \emph{inherent} to the task, or is DLoc simply
over-parameterized? If it is over-parameterized, how small a model can still localize accurately, and
what do we lose along the way?

\textbf{What we did.} We compressed DLoc into a family of lightweight ``student'' models
---MobileNetV2 (2.3\,M), a minimal TinyCNN (47\,K), a MobileNetV2 encoder with U-Net skip connections
(2.41\,M), and a state-space Mamba model (0.65\,M)--- and trained each one both \emph{standalone} and
\emph{distilled} from the DLoc teacher. We then evaluated all of them along five axes: in-distribution
accuracy, robustness to signal degradation, cross-session generalization, statistical reproducibility,
and computational efficiency.

\textbf{The headline result.} A \textbf{MobileNetV2-UNet with only 2.41\,M parameters reaches
0.640\,m median error}---\emph{better} than the 0.707\,m DLoc baseline---while using $6.8\times$ fewer
parameters and $56\times$ fewer floating-point operations. The cost of DLoc is therefore \emph{not}
inherent: a model an order of magnitude smaller can match or beat it.

\begin{keyidea}
The thesis is fundamentally an \textbf{efficiency} result wrapped around a \textbf{scientific} one.
The efficiency result: an order-of-magnitude-smaller model matches DLoc. The scientific result: the
two ways we made the small model good---\emph{U-Net skip connections} and \emph{knowledge
distillation}---improve different kinds of \emph{robustness}, and distillation only helps models above
a capacity threshold of about 2\,M parameters.
\end{keyidea}

\section*{The argument, as a chain}

Read the thesis as a single chain of claims; every chapter of this guide deepens one link:

\begin{enumerate}[leftmargin=2.2em]
  \item Localization can be framed as \emph{image-to-image translation} from CSI heatmaps to a
  location heatmap (Ch.~\ref{ch:foundations}).
  \item The building blocks that make a network small \emph{and} good are convolutions, skip
  connections, depthwise-separable convolutions, and U-Net skips (Ch.~\ref{ch:blocks}).
  \item \emph{Knowledge distillation} can transfer a big teacher's behavior into a small student
  (Ch.~\ref{ch:distill}); for our heatmap task it acts as ``spatial label smoothing.''
  \item We measure everything with a pixel-distance error and its \emph{distribution} ---median, P90,
  P99, mean, and the CDF (Ch.~\ref{ch:eval}).
  \item Each architecture then behaves in a specific, explainable way ---the over-parameterized
  teacher, the lossless MobileNetV2, the under-capacity TinyCNN, the teacher-beating U-Net, the
  mismatched Mamba (Ch.~\ref{ch:arch}).
  \item Comparing them yields six findings about compression, architecture, and distillation
  (Ch.~\ref{ch:findings}).
  \item Finally, anticipate the committee's questions (Ch.~\ref{ch:qa}).
\end{enumerate}

\section*{How to read it, depending on your goal}

\begin{itemize}[leftmargin=1.4em]
  \item \textbf{``I want to understand the whole thing.''} Read front to back. Each chapter assumes
  the previous one.
  \item \textbf{``I own the architecture questions.''} Read Ch.~\ref{ch:blocks} (concepts) then
  Ch.~\ref{ch:arch} (each model).
  \item \textbf{``I own distillation/robustness.''} Read Ch.~\ref{ch:distill} and Ch.~\ref{ch:eval},
  then findings 3--5 in Ch.~\ref{ch:findings}.
  \item \textbf{``I want to rehearse.''} Go straight to the Defense Q\&A (Ch.~\ref{ch:qa}) and the
  cheat-sheet in Appendix~\ref{app:glossary}; follow each answer's cross-reference back when a concept
  is shaky.
\end{itemize}

\section*{Concept-to-thesis map}

Table~\ref{tab:concept_map} points each concept to where it lives in the thesis, so this guide and the
thesis stay in lock-step.

\begin{table}[H]
\centering
\caption{Where each concept appears in the thesis.}
\label{tab:concept_map}
\small
\begin{tabular}{@{}lll@{}}
\toprule
\textbf{Concept} & \textbf{This guide} & \textbf{Thesis location} \\
\midrule
CSI, AoA/ToF, image-translation framing & Ch.~\ref{ch:foundations} & Sec.~3.1, docs/01--02 \\
Conv / encoder--decoder / skips         & Ch.~\ref{ch:blocks}      & Sec.~2.3 \\
Depthwise-separable, inverted residuals & Ch.~\ref{ch:blocks}      & Sec.~2.3.2 \\
U-Net skip connections                  & Ch.~\ref{ch:blocks}      & Sec.~2.3.3, 3.3.3 \\
State-space models / Mamba              & Ch.~\ref{ch:blocks}      & Sec.~2.3.4, 3.3.4 \\
Knowledge distillation                  & Ch.~\ref{ch:distill}     & Sec.~2.2, 3.4 \\
Error metric, percentiles, CDF          & Ch.~\ref{ch:eval}        & Sec.~3.6, 5.1 \\
Robustness corruptions                  & Ch.~\ref{ch:eval}        & Sec.~3.5, 4.4, 5.3 \\
Cross-session (LOSO / zero-shot)        & Ch.~\ref{ch:eval}        & Sec.~4.5, 5.4 \\
Per-model results \& reasons            & Ch.~\ref{ch:arch}        & Ch.~5 (all) \\
The six findings                        & Ch.~\ref{ch:findings}    & Ch.~5--6 \\
\bottomrule
\end{tabular}
\end{table}

A running convention: green bold marks a \best{best} result and red marks a \worst{worst} result, the
same as in the thesis tables.
